\documentclass[mathserif]{beamer}

\usepackage{listings}
\usepackage{showexpl}
\usepackage{setspace}
\usepackage{multirow}
\usepackage{dtklogos}
\usepackage[T1]{fontenc}

%To use handoutwithnotes: uncomments below
%\usepackage{handoutwithnotes}
%\pgfpagesuselayout{3 on 1 with notes}[a4paper,border shrink=5mm]
%\pgfpageslogicalpageoptions{1}{border code=\pgfusepath{stroke}}
%\pgfpageslogicalpageoptions{2}{border code=\pgfusepath{stroke}}
%\pgfpageslogicalpageoptions{3}{border code=\pgfusepath{stroke}}

\lstdefinestyle{latexsty}{
	language={[LaTeX]TeX},
    basicstyle=\small\ttfamily,
    breaklines=true,
    breakindent=0pt,
    backgroundcolor=\color{lightgray},
    numbers=none, numberstyle=\tiny, stepnumber=1, numbersep=5pt,
    commentstyle=\color{red},
    showstringspaces=false,
    keywordstyle=\color{blue}\bfseries,
    morekeywords={align,begin},
    tabsize=2,
    pos=b
}

\usetheme{default}
\useoutertheme{infolines}
\usecolortheme[RGB={166,5,20}]{structure}
%\setbeamertemplate{items}[circle]
\setbeamertemplate{blocks}[rounded][shadow=false]
\setbeamertemplate{navigation symbols}{}


\title{Introduction to \LaTeX}
\subtitle{A Method for Scientific Writing}
\author[Jiang Jiaming]{Jiang Jiaming, Group 2}
\institute[SKLSE]{State Key Laboratory of Software Engineering, Wuhan University}
\date[20/03/15]{March 20th, 2015}


\begin{document}

%--------------- slide -------------------
\begin{frame}[plain]
	\titlepage
\end{frame}
	

%--------------- slide -------------------
\begin{frame}{Introduction}

\vfill
\begin{itemize}
	\item What, Why, Where and How of \LaTeX
	\item Writing \LaTeX\ (Basics)
        \begin{itemize}
        	\item Comments
        	\item Commands
            \item Document Structure
        \end{itemize}
	\item Elements
        \begin{itemize}
        	\item Text Formatting
        	\item Tables and Lists
        	\item Graphics, Images \& Figures in \LaTeX
        	\item Maths
        \end{itemize}
\end{itemize}
\vfill
\end{frame}

%--------------- slide -------------------
\begin{frame}{Common Problems in Typesetting Documents with Word}

\begin{itemize}
	\item Poor typographic control.
	\item "Badness" is difficult to resolve.
	\item Image embedding (instead of external links) and limited image editing
options.
	\item Unexpected surprises when using external text.
	\item Unprofessional look.
\end{itemize}

\end{frame}


%--------------- slide -------------------
\begin{frame}{What is~\LaTeX?}

\TeX\ is software developed by Donald Knuth for typesetting documents.

\begin{itemize}
	\item low level markup language and compiler
	\item very powerful, but difficult to use
\end{itemize}
\vfill

\LaTeX\ is a collection of software built around~\TeX\ to make life easier
\begin{itemize}
	\item macros and scripts to convert~\LaTeX\ commands to~\TeX
	\item many packages for doing complex formatting/layouts
\end{itemize}
\vfill
\begin{center}
    \includegraphics[scale=0.22]{img/DonaldKnuth}
\end{center}
\end{frame}


%--------------- slide -------------------
\begin{frame}{Why~\LaTeX?}

\begin{itemize}
	\item Allows you to concentrate on content and structure, rather than layout and presentation
	\item Keeps formatting consistent throughout the document
	\item No design or typography knowledge required
	\item Excellent for writing mathematical expressions or equations
	\item Long and complex documents can be created easily
	\item Free!
\end{itemize}

\end{frame}

%--------------- slide -------------------
\begin{frame}{Why~\LaTeX?}

... It looks great.

\begin{center}
    \includegraphics[scale=0.32]{img/Latex_beautiful}
\end{center}

\end{frame}

%--------------- slide -------------------
\begin{frame}{Why not~\LaTeX?}

\begin{itemize}
	\item Can't easily see how document looks as we write it
	\item Need to learn \LaTeX\ before you can create anything
	\item Lose some control over formatting and layout
	\item \LaTeX\ is not the prettiest or simplest language out there
\end{itemize}

\end{frame}


%--------------- slide -------------------
\begin{frame}{Where do we get~\LaTeX?}

\LaTeX\ is freely available from a number of different sources. There are many different implementations and collections of packages online.

\begin{description}
	\item[Windows] \textbf{proTeXt} is a good solution for Windows, providing \textbf{MiKteX}.
	\item[Mac OSX] \textbf{MacTeX} provides \textbf{TeXLive}.
	\item[Linux] Most distributions will come with~\LaTeX\ ready installed but if not you can install \textbf{TeXLive}.
\end{description}
\vfill
Some websites like \textbf{Overleaf} provide online platforms for ''Real-time Collaborative Writing and Publishing''.
~\\
\textbf{CTeX} (Chinese TeX) is a nice choice for Chinese users.
\end{frame}



%--------------- slide -------------------
\begin{frame}{How do we use~\LaTeX?}

Creating documents with~\LaTeX\ is simple:
\vfill
\begin{enumerate}
	\item Write our document as plain text in a `\texttt{.tex}' file, using~\LaTeX\ commands to structure and format it 	\item Compile our `\texttt{.tex}' file to produce the output
\end{enumerate}
\vfill
    \begin{center}
		\includegraphics[height=0.45\textwidth]{img/HowLatex}
	\end{center}
\vfill
\end{frame}


%--------------- slide -------------------
\begin{frame}[fragile]
\frametitle{First (basic)~\LaTeX\ Example}
	\begin{LTXexample}[style=latexsty]
		\documentclass{article}
		\begin{document}
		    Hello World!
		\end{document}
	\end{LTXexample}
	
\end{frame}


%--------------- slide -------------------
\begin{frame}[fragile]
\frametitle{Writing \LaTeX\ (Basics) -- Comments, Commands}
\begin{itemize}
	\item Comments
\end{itemize}
The `\texttt{\%}' character is used to create comments in~\LaTeX. 

	\begin{LTXexample}[style=latexsty]
		%This is a comment and will not be shown.
		Here is some text in our file that will be shown. %but the rest of the line will not be.
	\end{LTXexample}
\vfill
~\\
\begin{itemize}
	\item Commands
\end{itemize}
\LaTeX\ commands have an effect on the text in the document. Some commands have additional arguments or optional parameters. The general syntax for a~\LaTeX\ command is:

\begin{lstlisting}[style=latexsty]
\commandname[opt1, opt2, ...]\{arg1\}\{arg2\}...}
\end{lstlisting}
\vfill


\end{frame}


%--------------- slide -------------------
\begin{frame}[fragile]
\frametitle{Writing \LaTeX\ (Basics) -- Document Structure}

Every~\LaTeX\ document must have a certain structure:
\vfill
\begin{lstlisting}[style=latexsty]
\documentclass{...}
\usepackage{...}
\begin{document}
...
\end{document}
\end{lstlisting}
\vfill

The \texttt{{\textbackslash{documentclass\{\ldots\}}}} command tells~\LaTeX\ which type of document we are creating, and how it should be set up and formatted.
\vfill
Often,the default set of commands available to~\LaTeX\ cannot solve all of our problems alone. You will need \texttt{{\textbackslash{usepackage\{\ldots\}}}} command to include extra packages.
\vfill
Documents can be divided into different parts, chapters or sections.
\end{frame}


%--------------- slide -------------------
\begin{frame}{fragile}
\frametitle{Document Structure Example}

\vfill
\LaTeX\ Template for \texttt{IEEEtran} Journals

\vfill
    \begin{center}
		\includegraphics[height=0.5\textwidth]{img/WholeExample}
	\end{center}
\vfill

\end{frame}

%--------------- slide -------------------
\begin{frame}[fragile]
\frametitle{Elements -- Text Formatting}
\vfill
\LaTeX\ has many text formatting options included by default that allow you to change text size, style and spacing.
\vfill
	\begin{LTXexample}[style=latexsty,pos=l]
		{\Large This text is large} \\
{\tiny This text is tiny}
	\end{LTXexample}

\vfill
	\begin{LTXexample}[style=latexsty,pos=l]
		\textsf{Sans Serif font} \\
\texttt{Teletype font} \\
\textit{Italic} \\
\textsc{Small Capitals} \\
\textbf{Bold} \\
\textrm{Roman font}
	\end{LTXexample}

\vfill
	\begin{LTXexample}[style=latexsty,pos=l]
		We can \emph{emphasise} a word in the sentence.
	\end{LTXexample}
\vfill
\end{frame}


%--------------- slide -------------------
\begin{frame}
\frametitle{Elements -- Lists}
\vfill
\LaTeX\ provides three different environments for defining lists:
\vfill
	\begin{description}
		\item[\texttt{itemize}] For basic bulleted lists
		\item[\texttt{enumerate}] For numbered lists
		\item[\texttt{description}] For lists with a label for each item
	\end{description}
\vfill
\end{frame}

%--------------- slide -------------------
\begin{frame}[fragile]
\frametitle{\texttt{itemize, enumerate, description}}


\vfill
	\begin{LTXexample}[style=latexsty,pos=l]
		\begin{itemize}
		    \item the first item
		    \item the second item
		    \item the third item
		\end{itemize}
	\end{LTXexample}

\vfill
	\begin{LTXexample}[style=latexsty,pos=l]
		\begin{enumerate}
		    \item the first item
		    \item the second item
		    \item the third item
		\end{enumerate}
	\end{LTXexample}

\vfill
	\begin{LTXexample}[style=latexsty,pos=l]
		\begin{description}
		    \item[First] the first item
		    \item[Second] the second item
		    \item[Third] the third item
		\end{description}
	\end{LTXexample}
\vfill

\end{frame}

%--------------- slide -------------------
\begin{frame}[fragile]
\frametitle{Elements -- Tables}
\vfill
The \texttt{tabular} environment allows us to create tables within our \LaTeX\ documents.
\vfill
\emph{Table} Example:
\vfill
	\begin{LTXexample}[style=latexsty]
		\begin{tabular}{ | l || c || r | }
		    \hline
		    some text & some text & this text \\
		    a text & other text & more text \\
		    another text & some text & some other text \\
		    \hline
		\end{tabular}
	\end{LTXexample}
\vfill
\end{frame}

%--------------- slide -------------------
\begin{frame}[fragile]
\frametitle{Elements -- Graphics, Images \& Figures in \LaTeX}
\vfill
Including graphics in our \LaTeX\ document is relatively simple.
\vfill
\begin{itemize}
	\item Insertion
\end{itemize}
The \texttt{{\textbackslash}includegraphics} command is used to insert an image file into your \LaTeX\ document.
\vfill
\begin{lstlisting}[style=latexsty]
	\includegraphics[opt1=val1, ..., ]{imagefile}
\end{lstlisting}
\vfill
\begin{itemize}
	\item Options
\end{itemize}
\begin{tabular}{r | l}
	\texttt{width=xx} & Specify the preferred width of the image \\
  	\texttt{height=xx} & Specify the preferred height of the image \\
    \texttt{scale=xx} & Scale the image by the given factor \\
    \texttt{angle=xx} & Rotate the image by the given number of degrees \\
	\texttt{page=x} & Select a particular page from a \texttt{.pdf} file\\
\end{tabular}
\vfill
\end{frame}

%--------------- slide -------------------
\begin{frame}[fragile]
\frametitle{\texttt{{\textbackslash}includegraphics} -- example}
\vfill
	\begin{LTXexample}[style=latexsty]
		\begin{center}
            \setlength\fboxrule{2pt}
            \fbox{\includegraphics[width=0.3\textwidth, angle=5]{img/background}}
		\end{center}
	\end{LTXexample}
\vfill
\end{frame}


%--------------- slide -------------------
\begin{frame}{Elements -- Maths}

\vfill
Typesetting maths is one of the major strengths of \LaTeX, and one of the main reasons for its popularity within the scientific community.
~\\~\\
The \texttt{amsmath} package is frequently used to add to the built in maths typesetting capabilities of \LaTeX -- all the discussion here will assume it has been used.
\vfill

\vfill
\emph{Besides}, there are many more packages that can be useful for this type of scientific work. Just some are: \texttt{IEEEtrantools}, \texttt{mathtools}, \texttt{mchem}.
\vfill

\end{frame}

%--------------- slide -------------------
\begin{frame}[fragile]
\frametitle{\texttt{Maths Examples}}

\vfill
    Inline maths
	\begin{LTXexample}[style=latexsty,pos=l]
		Inline maths is displayed within our text, such as \begin{math} a^3 + 6b \end{math}.
	\end{LTXexample}

\vfill
    displaymath(without numbering)
	\begin{LTXexample}[style=latexsty,pos=l]
		Display maths is displayed separately from the text.
		\begin{displaymath}
		    a^3 + 6b
		\end{displaymath}
	\end{LTXexample}

\vfill
    equation(with numbering)
	\begin{LTXexample}[style=latexsty,pos=l]
		Equation - separate, with numbering.
		\begin{equation}
		    a^3 + 6b
		\end{equation}
	\end{LTXexample}
\vfill

\end{frame}

%--------------- slide -------------------
\begin{frame}{fragile}
\frametitle{Special Documents}

\vfill
\begin{itemize}
	\item Presentation
	\item Poster
	\item CV and cover letter
	\item Teaching stuff
	\begin{center}
		\includegraphics[height=0.5\textwidth]{img/cv_template}
        \includegraphics[height=0.5\textwidth]{img/poster}
	\end{center}
\end{itemize}
\vfill
\end{frame}

%--------------- slide -------------------
\begin{frame}{fragile}
\frametitle{Questions?}

\vfill
    \begin{center}
		\includegraphics[height=0.5\textwidth]{img/QA}
	\end{center}
\vfill
\end{frame}

\end{document}
